\chapter{Appendices}
\label{ch:appendices}

\section{Appendix A --- Glossary of Terms}
\label{app:glossary}

\begin{description}
  \item[Wait-state advertising] The category Devorix operates in: placing a
        sponsor message in the brief ``thinking'' / idle status moment of an AI
        coding tool, when the developer is waiting for the agent to respond.
        The defining feature is that the attention is captive and full ---
        unlike a sidebar the user scrolls past.

  \item[Status line / status-bar overlay] The thin text slot an AI coding tool
        renders during a thinking state (and the secondary VS~Code status-bar
        item). Devorix injects its sponsor line here via the configurable
        \texttt{statusLine} mechanism, \emph{non-patching} (see below).

  \item[Non-patching client] An integration that injects the sponsor line
        through a supported configuration hook and \emph{never} modifies the
        host tool's rendering layer, never reads code or prompts, and never
        touches the host tool's OAuth tokens. This is the basis of Devorix's
        platform-policy and ``we are not Kickbacks'' compliance posture.

  \item[Impression] A single billable display of a sponsor creative. A display
        counts as billable only when the creative has shown continuously for at
        least the billable impression length (5~seconds for v1) during an
        active session, and only after it passes server-side validity
        filtering.

  \item[Billable impression length] The minimum continuous display time before
        a display counts as a billable impression. Locked at \textbf{5~seconds}
        for the v1 build; 10~seconds retained as a documented, data-gated
        option.

  \item[Click] A counted event when a developer opens the optional click URL
        attached to a creative.

  \item[Fill rate] The percentage of available impression slots that are
        served a \emph{paid} (non-house) advertiser creative rather than a
        house ad. The single most sensitive variable in the financial model and
        a direct proxy for whether real advertiser demand exists; today
        $\sim$0\% (no advertisers).

  \item[House ad] A Devorix-controlled promotional creative served when no paid
        advertiser creative is available (e.g.\ at launch, before any
        advertiser exists). Generates a billable-shaped event at ₹0
        advertiser cost, used to validate the loop.

  \item[Tier 0 / Tier 1 / Tier 2] The three rungs of the locked hybrid pricing
        ladder. \textbf{Tier~0}: flat monthly sponsorships (₹10{,}000--
        15{,}000/mo) plus house ads, until $\sim$5 advertisers.
        \textbf{Tier~1}: a fixed-price CPM rate card (\$2--3 India / \$5--8
        Global), once 5--10 advertisers are signed. \textbf{Tier~2}: a
        real-time second-price auction, once 5{,}000+ DAU.

  \item[CPM] Cost per mille --- price per 1{,}000 impressions. The pricing
        basis for Tier~1 and above. The only independently audited real
        developer-advertising comp is Carbon Ads at \$1.60 CPM.

  \item[eCPM] Effective CPM --- revenue normalized to a per-1{,}000-impression
        figure, used to compare flat-fee Tier~0 deals against a CPM basis.

  \item[Block] A Tier~1 unit of inventory equal to 1{,}000 impressions; the
        unit an advertiser buys against the rate card.

  \item[DAU / WAU] Daily / Weekly Active Users --- developers with the client
        enabled and at least one session in the period. The 5{,}000+ \textbf{DAU}
        threshold triggers Tier~2.

  \item[Payout threshold] The accrued payable balance (locked at
        \textbf{₹300}) at which an automatic UPI payout is triggered.

  \item[Payable, not stored value] The regulatory framing of the developer's
        accrued balance: a debt Devorix owes, settled to a bank/UPI account on
        the threshold, with no load/spend/hold capability --- deliberately
        outside RBI's Prepaid Payment Instrument (PPI) definition.

  \item[UPI / VPA] Unified Payments Interface, India's real-time payment rail;
        a VPA (Virtual Payment Address, e.g.\ \texttt{name@bank}) is the
        UPI handle a payout is sent to.

  \item[RazorpayX] The payout provider; its Payouts API executes automated
        UPI/IMPS/NEFT/RTGS transfers with 24/7 instant UPI settlement.

  \item[Double-entry ledger] The money source of truth: every economic event
        posts as a balanced pair of entries, and a developer's visible balance
        is always a derived read (\,$\sum$ credits $-\,\sum$ debits\,), never a
        separately stored mutable field.

  \item[Raw / billable event separation] The architectural rule that raw
        impression events are stored separately from the fraud-filtered
        \texttt{billable\_events} that alone drive advertiser billing and
        developer payouts.

  \item[Atomic network] (Andrew Chen, \emph{The Cold Start Problem}) The
        smallest dense network that can stand on its own; Devorix's
        India-first, Ring-1-first plan is structurally an atomic-network
        strategy.

  \item[Three-ring niche] Devorix's advertiser targeting:
        \textbf{Ring~1} (AI-infra/dev-tool sellers to AI-building developers),
        \textbf{Ring~2} (India dev-economy sellers), \textbf{Ring~3} (global
        dev-tool brands, Phase~2 only).

  \item[Go/no-go gate] The locked month-4 test that governs continuation: at
        least two advertisers paid \emph{and} renewed, week-4 retention holds,
        and the founder still wants to run it.

  \item[Sweat equity] The founder's labour, modelled at ₹0 salary, so all
        margins are not fully loaded and exclude founder opportunity cost.
\end{description}

\section{Appendix B --- Locked Decisions (verbatim)}
\label{app:locked-decisions}

The table below reproduces \texttt{MASTER-PLAN.md} \S2 in full, with all
numbers verbatim. No decision in this table may be altered or hedged anywhere
in this document; they are stated as final facts.

{\small
\begin{longtable}{p{0.4cm} p{5.1cm} p{2.4cm} p{4.6cm}}
\caption{Locked decisions (\texttt{MASTER-PLAN.md} \S2, verbatim).}
\label{tab:locked-decisions} \\
\toprule
\textbf{\#} & \textbf{Decision} & \textbf{Status} & \textbf{Notes} \\
\midrule
\endfirsthead
\toprule
\textbf{\#} & \textbf{Decision} & \textbf{Status} & \textbf{Notes} \\
\midrule
\endhead
1 & Model = ads-to-devs (Kickbacks-style) & Decided & --- \\
\addlinespace
2 & First tool = Claude Code; client = non-patching/status-bar & Decided &
Enterprise-safe, avoids Kickbacks' ToS risk \\
\addlinespace
3 & Payout rail = RazorpayX UPI, accrue + batch & Decided & Re-verified
2026-06-28: RazorpayX Payouts API confirmed to support automated
UPI/IMPS/NEFT/RTGS, 24/7 instant settlement \\
\addlinespace
4 & Build = solo, founder-built MVP & Decided & --- \\
\addlinespace
5 & Public framing = ``sponsorship footer'' & Decided & From AI Studio
prototype \\
\addlinespace
6 & Adopt 3 Gemini prototype tools & Decided & Ad-Copy Generator (P1),
Wait-State Coach (P1), Workspace Analyzer (P2); core ad-serving/ledger must not
depend on Gemini \\
\addlinespace
7 & Entity/brand = ``Devorix'' / ``DEVORIX.'' & \textbf{DECIDED FINAL} & Renamed
from ``Async Returns'' / ``ASYNC.'' on 2026-06-28 per founder's explicit call.
Domain placeholder used across docs/code: devorix.ai --- not yet verified as
registered/available. \\
\addlinespace
8 & Budget/runway $\leq$ ₹1~lakh, low burn, solo & Founder constraint &
--- \\
\addlinespace
9 & Do NOT incorporate yet --- sole proprietor & Decided & Re-verified
2026-06-28, see \S7.1 \\
\addlinespace
10 & No GST registration yet (below ₹20L/yr threshold) & Decided &
Re-verified 2026-06-28, see \S7.1 \\
\addlinespace
11 & \textbf{Dev rev share = 50 / 50} & \textbf{DECIDED FINAL} & Closes the
50-vs-60-70\% open question. Overrides FINANCIAL\_MODEL's 40/60 and 01/03's
``60-70\% as install-war lever'' idea. \\
\addlinespace
12 & \textbf{Geographic sequencing: India first, Global = explicit Phase 2} &
\textbf{DECIDED FINAL} & Gated on the go/no-go criteria in \S8.3. PRD's
``non-India = non-goal'' is correct for now, not permanent. \\
\addlinespace
13 & \textbf{Payout threshold = ₹300} & \textbf{DECIDED FINAL} & Overrides
₹500 (05) and ₹250 (PRICING\_ARCHITECTURE) --- founder's call, splits
the difference \\
\addlinespace
14 & \textbf{Pricing model = hybrid ladder} --- Tier~0 (house ads +
opportunistic flat-fee deals, until $\sim$5 advertisers) $\rightarrow$ Tier~1
(real fixed-price CPM rate card, once 5--10 advertisers) $\rightarrow$ Tier~2
(real-time auction, once 5{,}000+ DAU) & \textbf{DECIDED FINAL} & Full detail
in \S5 \\
\addlinespace
15 & \textbf{Billable impression length = 5 seconds for v1 build}; 10 seconds
kept as a documented, data-gated option & \textbf{DECIDED FINAL for now} &
Founder asked to compare, not force a permanent choice --- see \S5.3 \\
\addlinespace
16 & \textbf{USD$\rightarrow$INR = ₹94} & \textbf{RE-VERIFIED 2026-06-28} &
Live mid-market $\sim$₹94.36--94.44. Replaces ₹86 (00), ₹83
(PRICING\_ARCHITECTURE, FINANCIAL\_MODEL); confirms ₹94.5 (05) was
right. \\
\addlinespace
17 & \textbf{Warm advertiser intros at sprint start = 0} (fully cold outreach) &
Founder constraint, newly confirmed & Adjusts sprint math, see \S6.3 \\
\addlinespace
18 & \textbf{Tier-1 rate-card CPM (once activated) = $\sim$\$2--3 India /
\$5--8 Global} & \textbf{DECIDED FINAL (planning number)} & Replaces \$8/\$15
(PRICING\_ARCHITECTURE) and \$5$\rightarrow$\$10-12 (FINANCIAL\_MODEL) ---
anchored to the one independently-audited real comp, not 4--10x it. Full
reasoning in \S5.4. Tagged \emph{Assumption} --- revisit with real fill
data. \\
\bottomrule
\end{longtable}
}

\section{Appendix C --- Methodology Note on Evidence Labelling}
\label{app:methodology}

Throughout this document, every non-trivial factual claim --- and especially
every number --- carries one of three labels, used consistently and never
blurred together through synthesis:

\begin{description}
  \item[Verified fact] Named, on-record, or independently multi-sourced. These
        are claims traceable to a primary source (e.g.\ a company's own
        disclosure, a regulator's documentation, or two or more independent
        reports). Example: ``India has 27~million developers on GitHub (2026)''
        --- attributed on record to GitHub's COO and corroborated across
        multiple outlets.

  \item[Industry estimate] Published by research firms or third parties,
        directionally consistent but without a single authoritative number, or
        self-reported. These are credible and useful as range/context but should
        not be treated as precise. Example: the AI code tools market at
        ``$\sim$\$9--9.5B (2026)'' across firms whose category boundaries
        disagree.

  \item[Internal assumption] Devorix's own model output, reasoning, or
        extrapolation --- not an externally researched fact. All figures past
        the locked, audited month-4 plan fall here. Example: the 12--24 month
        revenue scenarios, which are planner extrapolation, not workbook
        outputs.
\end{description}

Two conventions follow from this labelling. First, \textbf{the locked decisions
(Appendix~\ref{app:locked-decisions}) and the audited months-1--4 plan are the
only load-bearing inputs to the funding ask}; everything labelled Internal
assumption (the TAM ceiling, the 12--24 month curves, the agentic-economy
context) is optionality and direction, presented to explain \emph{why the bet
is worth taking if the near-term loop works}, never as something derived from
the financial workbooks. Second, where two source inputs disagreed, the
contradiction is \emph{resolved explicitly in the text} (e.g.\ the 18\%
workplace-usage vs.\ 28\%/24\% primary-tool-split adoption figures, reconciled
as measuring different populations) rather than silently dropped --- with the
more conservative, better-sourced figure adopted as the planning anchor and the
favourable figure carried forward only as a directional signal.

\section{Appendix D --- Note on the Multi-Agent Research Process}
\label{app:research-process}

In the interest of transparency about how this PRD was produced, this appendix
documents the research pipeline behind it.

The document was synthesized through a deliberate multi-agent workflow rather
than written in a single pass. The pipeline had three roles:

\begin{enumerate}
  \item \textbf{Explorer agents} performed broad, read-only fan-out research ---
        gathering named comparables, competitor launch mechanics, fund theses,
        regulatory specifics, and market figures from primary sources --- and
        returned findings with source attribution and confidence labels. Their
        job was to locate and verify evidence, not to draw conclusions.

  \item \textbf{Planner agents} structured that evidence into chapter-level
        synthesis: building the 24-month roadmap contingency tree, the moats
        verdict, the financial scenario extension, the persona set, and the
        risk register, while keeping the explorer's confidence labels intact and
        flagging open questions rather than papering over them.

  \item \textbf{An orchestrating synthesis pass} reconciled the planner outputs
        against the single source of truth (\texttt{MASTER-PLAN.md}), resolved
        every cross-document contradiction explicitly, enforced the locked
        decisions verbatim, and produced the final chapter text.
\end{enumerate}

The single source of truth throughout was \texttt{MASTER-PLAN.md}: where any
specialist input disagreed with it, the master plan won, and no locked decision
in Appendix~\ref{app:locked-decisions} was reopened at any stage. This process
is disclosed here so a reader can weigh the document's claims knowing how they
were assembled --- in particular, that the evidence labelling in
Appendix~\ref{app:methodology} was applied at the point of research collection
and preserved through synthesis, not retrofitted at the end.
